\section{Alternating quotients of arbitrary tensor-permutation images}
\label{sec:tensor-permutation-alternating-quotient}

\paragraph{Status and provenance.}
D-TD45, reserved question Q2fa, is a separate extension of D-TD15
by \texttt{/root/g34\_eightfolds}. The new step applies root's
accepted D-TD40(a) to the actual transitive permutation image,
and combines its uniform alternative with the complete-lift trace
identity of D-TD14. The latter is an imprimitive derivation.
The consequence for coordinate quotients uses the accepted D-TD17--18
and D-TD43. No new classification theorem or numerical experiment
is introduced. Complete author whole-proof and source-interface
readings pass. Independent and root whole-proof reviews also pass.

\subsection*{D-TD45. The tensor-orbit quotient obstruction}

Let a finite group $J$ act on a complex space with an actual tensor
decomposition
\[
 V=\left(\bigotimes_{\omega\in\Omega}W_\omega\right)\otimes R,
 \qquad |\Omega|=t,\qquad \dim W_\omega=a\ge2,
 \qquad e=\dim V=a^t\dim R.
 \tag{TPQ.1}
\]
Assume $J$ permutes the displayed embedded full matrix algebras,
with actual transitive image
\[
 \pi:J\twoheadrightarrow\Pi\le\operatorname{Sym}(\Omega),
 \qquad \varphi:\Pi\twoheadrightarrow T\in\{A_n,S_n\},
 \qquad x=\varphi\pi.
 \tag{TPQ.2}
\]
The entire external partner $R$ is arbitrary. Its full matrix algebra
is the commutant of the product of the displayed factor algebras,
and is preserved. Thus the original operators have the weighted
tensor-permutation form used in D-TD14, with all scalar ambiguities
and the external factor retained. No factor $W_\omega$ is asserted
to be an individual representation of $J$.

\begin{theorem}[D-TD45]
If an actual connected $J$-cover of the sphere has original inertia
generators satisfying
\[
              \sum_i\operatorname{rank}(1-V(g_i))=2e+6,
 \tag{TPQ.3}
\]
then (TPQ.1)--(TPQ.2) require $n<10^6$.
The entire permutation image $\Pi$ need not be alternating,
symmetric, primitive or quasiprimitive. The conclusion uses the
specified actual quotient in (TPQ.2), not an abstract composition
factor, and assumes no minimum number of projected branch entries.
\end{theorem}

\begin{proof}
Suppose $n\ge10^6$. Faithfulness of the permutation image gives
\[
               t!\ge|\Pi|\ge|T|\ge n!/2,
                    \qquad t\ge n,
                    \qquad e\ge2^t\ge2^n.
 \tag{TPQ.4}
\]
Indeed $t<n$ would imply $t!\le(n-1)!<n!/2$.
Apply D-TD40(a) to the actual transitive pair $(\Pi,\varphi)$.
It supplies one of the following alternatives for the whole pair:
\[
 \begin{split}
 &u\mid t,\quad n\le u\le t,\quad
 \operatorname{mov}_\Omega(\pi(g))
       \ge\frac{t}{u}\frac{\sqrt u-1}{2}
       \quad\text{whenever }x(g)\ne1;\quad\text{or}\\
 &n\mid t,\quad
 \operatorname{mov}_\Omega(\pi(g))
       \ge\frac{t}{n}\operatorname{mov}_{[n]}(x(g))
       \quad\text{for every }g\in J.
 \end{split}
 \tag{TPQ.5}
\]
The divisor $u$ and the alternative do not change with $g$.
The proof of that theorem follows actual block-kernel quotients
through a finite transitive-action induction; no primitive image
is substituted for $\Pi$ here.

For every original element $g$, and every integer $j$, the exact
weighted-tensor trace bound D-TD14 gives
\[
 \frac{|\operatorname{tr}V(g^j)|}{e}
    \le a^{-\operatorname{ind}_\Omega(\pi(g)^j)}
    \le2^{-\operatorname{mov}_\Omega(\pi(g)^j)/2}.
 \tag{TPQ.6}
\]
This is valid for the full original finite order. Its proof
reciprocally normalizes the two tensor operators to that order,
uses the actual cycle products, and includes their scalar moduli.
It does not require a tensor-product metric or claim that a power
with trivial factor permutation is scalar. All such powers retain
the normalized trace bound one.

\paragraph{The uniform-support alternative.}
In the first case of (TPQ.5), every $x(g)\ne1$ has
\[
 \operatorname{ind}_\Omega(\pi(g))
    \ge\frac{t}{u}\frac{\sqrt u-1}{4}
    \ge\frac{999}{4}>7,
 \qquad
 \frac{|\operatorname{tr}V(g)|}{e}<\frac1{128}.
 \tag{TPQ.7}
\]
Let $L$ be the original order of $g$ and $o$ the order of $x(g)$.
Exactly $L/o$ of the powers numbered $0,\ldots,L-1$ have trivial
natural image. Fourier inversion over all $L$ powers, with those
powers bounded by one and the others by $1/128$, gives for every
target eigenvalue, and hence for actual moving rank,
\[
          \frac{\operatorname{rank}(1-V(g))}{e}
                  \ge\frac{127}{128}(1-1/o).
 \tag{TPQ.8}
\]
The actual connected $T$-quotient cover has genus at least two:
$T$ is nonsolvable and its order exceeds sixty. The full actual-cover
argument of D-TD02 therefore supplies the orbifold sum at least
$85/42$. Summing (TPQ.8) and using (TPQ.3) yields
\[
 2+6/e\ge\frac{127\cdot85}{128\cdot42}
          =2+\frac{43}{5376},
 \qquad e\le\frac{32256}{43}<751.
 \tag{TPQ.9}
\]
This contradicts $e\ge2^n\ge2^{10}>751$.
Identity quotient entries contribute zero to the orbifold sum;
they are not removed from the original rank budget.

\paragraph{The natural-support alternative.}
In the second case of (TPQ.5), put
$k=\operatorname{mov}_{[n]}(x(g))$ and $f=n-k$.
Equations (TPQ.4)--(TPQ.6) give
\[
 \frac{|\operatorname{tr}V(g)|}{e}
           \le2^{-tk/(2n)}\le2^{-k/2}.
 \tag{TPQ.10}
\]
If $x(g)\ne1$ then $k\ge2$, so this is at most $1/2<7/8$.
The numerical comparison in D-TD14, (TO.8)--(TO.11), applies to
the present natural degree $n$. For completeness its full argument
is short. For every integer $n\ge128$, $2^n>n^{12}$: it holds
at $128$, and $(1+1/n)^{12}<(1-12/n)^{-1}\le32/29<2$
preserves it by induction. For $1\le f\le n$, the function
\[
       \frac{n-f}{2}\log2+3\log(f/n)
 \tag{TPQ.11}
\]
is concave. It is zero at $f=n$ and positive at $f=1$, since
$2^{n-1}>n^6$. It is therefore nonnegative throughout this
interval. For $f=0$, the same numerical estimate gives
$2^{-n/2}<n^{-6}$. Consequently every original element, including
every full original power, satisfies
\[
 \frac{|\operatorname{tr}V(g)|}{e}
       \le\left(\frac{\max(f,1)}{n}\right)^3,
 \qquad
 \frac{|\operatorname{tr}V(g)|}{e}\le\frac12
          \quad\text{when }x(g)\ne1,
 \qquad e\ge2^n>n^3.
 \tag{TPQ.12}
\]
These are precisely the cubic and uniform trace hypotheses of
HD-FT27 at $n\ge10^6$, with the entire partner already included
in $V$. That theorem contradicts (TPQ.3).

If the initial coefficient representation is not faithful, use
the actual connected cover $C/\ker V$. An operator acting as the
identity on $V$ centralizes every displayed full matrix algebra,
so it has trivial factor permutation and trivial image under $x$.
Thus $\ker V\subseteq\ker\pi\subseteq\ker x$. Both actual
quotient maps survive, the rank sum is unchanged, and identity
coefficient inertia images contribute zero. This is the exact
faithful-image passage required by HD-FT27. No new geometric
source or quotient-order Fourier period is being presumed.
\end{proof}

\subsection*{A coordinate quotient without a tensor-count premise}

\begin{corollary}[D-TD45, coordinate consequence]
In an actual source (TPQ.3), let a stable whole tensor factor $B$
have a finite primitive projective image
$P=\operatorname{im}(J\to\operatorname{PGL}(B))$ and an actual
surjection $\psi:P\twoheadrightarrow T\in\{A_n,S_n\}$,
with $n\ge2^{30}$. Fix a complete D-TD12 decomposition of $B$
into its original coarse extraction spaces. Then $\psi$ factors
through a coarse coordinate $P_j$ for which
\[
             N_j\simeq A_n,
             \qquad P_j\mathop{\longrightarrow}^{\sim}T.
 \tag{TPQ.13}
\]
Here $N_j$ is the actual nonabelian minimal normal extraction group,
with just one simple component. No bound on the number of components
of any other coordinate, and no maximality of $n$, is assumed.
This does not classify the irreducible projective module of $P_j$.
\end{corollary}

\begin{proof}
The complete coarse images give the actual subdirect embedding
of D-TD17. Its coordinate-kernel lemma makes $\psi$ factor through
some actual $P_j$. Each coordinate is stable under the original
group and has its whole external complement retained in $V$.

If it is a Heisenberg coordinate, the proof of D-TD18 supplies
its actual irreducible prime-field symplectic image and a genuine
surjection from that image onto $T$. It satisfies D-TD43 in the
same original source, which is impossible at $n\ge2^{30}$.

Otherwise its actual normal extraction group is $N=S^t$ for a
nonabelian simple $S$, with transitive component action. The
proof of D-TD18, (TDC.7)--(TDC.8), gives
$N\le P_j\le\operatorname{Aut}(N)$ and trivial projective
centralizer of $N$. Write $\alpha:P_j\twoheadrightarrow T$
for the selected actual quotient.

If $\alpha(N)\ne1$, its perfect normal image is $A_n$.
The simple-factor images are normal in $A_n$ and are all nontrivial
by component transitivity. Two such images would commute, so $t=1$
and $S\simeq A_n$. The restriction $\alpha|_N$ is injective.
Every element of $\ker\alpha$ then centralizes $N$, since its
commutators with $N$ lie in $N\cap\ker\alpha=1$. The trivial
projective centralizer gives $\ker\alpha=1$, proving (TPQ.13).

If $\alpha(N)=1$, the kernel of the actual component-permutation
map $P_j/N\to\Pi\le S_t$ embeds in $\operatorname{Out}(S)^t$.
The exact Schreier input already used by D-TD18 makes that kernel
solvable. Its image in $T$ is a solvable normal subgroup and is
therefore trivial. Thus the actual transitive component image
$\Pi$ has a genuine surjection onto $T$.

The coordinate space $B_j$ is the tensor product of the represented
quasisimple component modules. They have equal dimensions $a\ge2$,
their full matrix algebras are actually permuted with image $\Pi$,
and the original map $J\to P_j\to\Pi\to T$ is onto. All other
coarse factors and the original external partner form the entire
factor $R$ in (TPQ.1). The first theorem therefore excludes this
case at $n\ge2^{30}>10^6$. Unlike the earlier proof of D-TD18,
this step requires no inequality $t<n$.
\end{proof}

\paragraph{Strict scope and remaining boundary.}
D-TD15 required the entire tensor-permutation image to be $A_t$
or $S_t$. The present theorem permits an arbitrary transitive
image and an actual alternating or symmetric quotient of possibly
smaller degree. The support comparison is the proved D-TD40(a)
alternative, not an identification of those two permutation actions.
The coordinate consequence removes the tensor-count premise while
retaining a singleton projective alternating module of unrestricted
character type. In particular it does not declare a near-maximal
layer ordinary standard, change a coefficient field or polarization,
iterate a rank-nine source on a positive-cohomology partner, or bound
all remaining source dimensions.
